\section*{Chapitre \ref{ch:pl:native-extensions} --- Extensions natives}

\begin{solution}{exo:pl:native-extensions:1}
Par élément via pybind11~: $118 \times 2\times10^8$~ns $= 23.6$~s~; en Python~: 12.0~s~; un seul appel sur tout le tableau à 1.4~ns~: 0.28~s.
\end{solution}

\begin{solution}{exo:pl:native-extensions:2}
\texttt{ctypes} convertit chaque argument à l'exécution selon les types déclarés, par un mécanisme générique, et construit dynamiquement un cadre d'appel C
Les conversions de pybind11 sont compilées dans la liaison pour la signature exacte. Les instructions machine du noyau sont les
mêmes~; le travail qui les entoure, lui, ne l'est pas.
\end{solution}

\begin{solution}{exo:pl:native-extensions:3}
Il convertit le tableau en un nouveau tableau float64 et transmet la copie. Pour \texttt{ewma\_into}, le noyau écrit alors dans la copie, qui est
jetée, et le tableau de l'appelant reste inchangé~: un résultat faux sans aucun signe d'erreur.
\end{solution}

\begin{solution}{exo:pl:native-extensions:4}
Pour $n > c_0/(p - c_1) = 118/58.6 \approx 2.0$~: à partir de trois éléments, un seul appel sur tout le tableau bat la boucle Python.
\end{solution}

\begin{solution}{exo:pl:native-extensions:5}
La recherche de numpy consiste en une recherche dichotomique pour chaque instant de gauche, en $O(\log m)$ étapes chacune, et chaque étape touche une partie différente d'un tableau plus grand
(davantage de défauts de cache à mesure qu'il grandit)~; la fusion parcourt une seule fois, vers l'avant, les deux tableaux triés, pour un coût constant par élément avec un
accès mémoire séquentiel.
\end{solution}

\begin{solution}{exo:pl:native-extensions:6}
Après l'appel, Python peut libérer ou réallouer le tableau (ramasse-miettes, redimensionnement, remplacement), laissant au noyau un
pointeur pendant~: des lectures de mémoire libérée ou des écritures dans celle d'un autre. La règle~: un noyau lit et écrit les tampons seulement pendant l'appel et ne conserve
aucun pointeur ensuite~; si les données doivent persister, l'appelant garde le tableau en vie et le retransmet.
\end{solution}

\begin{solution}{exo:pl:native-extensions:7}
Les quatre s'accordent sur chacun des 100\,000 indices~: $-1$ pour l'instant précédant le premier instant de droite, et 49\,999 (le dernier indice) pour
l'instant qui suit le dernier.
\end{solution}

\begin{solution}{exo:pl:native-extensions:8}
Convertir une liste de flottants Python en vecteur puis reconvertir le résultat en liste revient à copier et à créer un objet par élément, le coût que le code natif devait éviter~;
une mesure sur dix nombres ne mesure que le coût fixe de l'appel, et non le coût par élément qui compte sur des données réelles. Il faut
prendre et rendre des tableaux numpy par le \omterm{def:pl:native-extensions:buffer}{protocole de tampon}, et mesurer sur des longueurs réalistes.
\end{solution}

\begin{solution}{pb:pl:native-extensions:1}
\begin{enumerate}
\item Elle localise la fonction, vérifie et convertit chaque argument à partir d'un objet Python, peut libérer puis reprendre le verrou, convertit le
résultat en retour et vérifie les erreurs.
\item Environ 60~ns (une fonction Python), 118~ns (C++ par pybind11) et 337~ns (Rust par \texttt{ctypes}), boucle comprise.
\item La boucle Python qui l'entoure subsiste, et chaque appel ajoute des conversions dont le coût dépasse celui de l'étape elle-même.
\item D'autres threads Python s'exécutent pendant que le code natif calcule~; cela n'est sûr que si le code natif ne touche aucun objet Python tant que le verrou
est libéré.
\item La convention d'appel et de disposition au niveau machine~; celle du C est simple, stable et reconnue par l'\omterm{def:pl:native-extensions:ffi}{interface de fonctions étrangères} (FFI) de tous les langages.
\item Le type, la contiguïté et le dtype de l'argument~: selon ces critères, il convertit ou refuse.
\item Que chaque tableau soit un tableau numpy du bon dtype et contigu en C, et transmettre sa longueur.
\item Les deux le refusent~: pybind11 parce que l'argument est \texttt{c\_style} sans \texttt{forcecast} ni \texttt{noconvert}~; l'enveloppe
parce qu'elle vérifie la contiguïté.
\item Elle copierait sans prévenir~: une sortie écrite dans une copie est perdue, et la copie d'un grand tableau coûte ce que l'extension devait
faire gagner.
\item Vérifications et conversions d'arguments générées pour Rust, comme celles de pybind11 pour C++, et types Rust pour les objets Python~; en dessous, la même
interface C.
\item Environ 1.4~ns (C++), 1.5~ns (Rust) et 3.3~ns (le filtre de numpy).
\item La routine de filtrage a un coût fixe de plusieurs microsecondes par appel, supérieur à celui des liaisons, et une boucle interne plus lente.
\item Environ 21~ns (numpy), 6~ns (C++) et 6~ns (Rust).
\item À partir d'environ mille éléments.
\item La recherche~: sa version native a un meilleur algorithme (une fusion au lieu d'une recherche dichotomique par élément), et pas seulement une boucle plus rapide.
\item \textbf{Résultat nommé.} Un appel coûte environ 60~ns comme fonction Python, 118~ns avec pybind11 et 337~ns avec \texttt{ctypes}~; sur
des tableaux entiers, les deux versions natives battent numpy à chaque longueur mesurée, et la recherche Rust à partir d'environ mille éléments~; déplacer la boucle plutôt que le corps
transforme une journée de 24~s d'appels élément par élément en 0.3~s.
\item Environ 12~s face à la boucle Python (12.0~s à 0.28~s), 23~s face aux appels natifs élément par élément, et 0.4~s face à numpy.
\item Quand le calcul est déjà une opération sur tableaux d'une bibliothèque, quand les tableaux sont courts et les appels peu nombreux, ou quand le gain ne
compense pas la construction, les tests et la distribution du code compilé sur chaque plateforme qu'utilisent les chercheurs.
\item Le noyau ne conserve aucun pointeur après l'appel et ne libère rien qu'il n'a pas alloué~; l'appelant garde ses tableaux vivants et inchangés
pendant l'appel.
\item Autour de tableaux entiers~: les données franchissent la frontière une seule fois, les boucles restent natives.
\end{enumerate}
\end{solution}

\begin{solution}{iq:pl:native-extensions:1}
Elle était appelée une fois par élément depuis la boucle Python~: chaque appel paie un coût de conversion et d'aiguillage supérieur à l'arithmétique, et la
boucle restait en Python. Déplacer la boucle dans la fonction native et transmettre des tableaux entiers.
\iqlookfor{Coût par appel contre coût par élément.}
\end{solution}

\begin{solution}{iq:pl:native-extensions:2}
Par le \omterm{def:pl:native-extensions:buffer}{protocole de tampon}~: transmettre le pointeur, le dtype et la forme du tableau (\texttt{array\_t} de pybind11, ou pointeurs \texttt{ctypes} issus de
\texttt{arr.ctypes}). Erreurs typiques~: un dtype ou des strides incorrects (lus comme s'ils étaient contigus), des copies silencieuses (forcecast), l'écriture dans un tableau en lecture seule, et des pointeurs
conservés après l'appel.
\iqlookfor{Vérification de la disposition et propriété des données.}
\end{solution}

\begin{solution}{iq:pl:native-extensions:3}
pybind11~: liaisons C++ compilées avec contrôles de type, conversions et gestion du verrou, la norme pour C++. \texttt{ctypes}~: pas de compilation d'une
liaison, fonctionne avec toute interface C, mais ne vérifie rien et coûte davantage par appel. Un générateur Rust (PyO3)~: le confort d'usage de pybind11 pour Rust~;
l'interface C avec \texttt{ctypes} fonctionne sans lui.
\iqlookfor{Sécurité et coût par appel contre commodité et complexité de compilation.}
\end{solution}

\begin{solution}{iq:pl:native-extensions:4}
Pour que d'autres threads Python s'exécutent pendant qu'il calcule. Il ne doit toucher aucun objet Python, compteurs de références compris, lever
d'exceptions Python, ni rappeler Python tant qu'il n'a pas repris le verrou.
\iqlookfor{Aucune API Python sans le verrou.}
\end{solution}

\begin{solution}{iq:pl:native-extensions:5}
Noyaux choisis par profilage (boucles que les tableaux ne savent pas exprimer, fusions, machines à états)~; écrits en C++ ou en Rust sur de simples pointeurs, testés
isolément~; liés une fois par tableau avec pybind11 (ou une interface C), refusant les dispositions qu'ils ne peuvent lire sur place et libérant le verrou~; compilés dans
le pipeline d'intégration continue pour chaque plateforme et chaque version de Python prise en charge, distribués sous forme de wheels~; testés face à une
référence Python sur les mêmes données.
\iqlookfor{Sélection par la mesure, séparation du noyau et de la liaison, constructions reproductibles.}
\end{solution}
